% Source PDF page 37; manual page 2-8.
\subsection{Local Geocentric Horizon Coordinate System}\label{sec:local-geocentric}

The local geocentric horizon coordinate system is defined to make geocentric
flight-path angles easy to compute. Its origin is at the vehicle's center of mass,
located by the position vector $\vec{r}$. The $z$ axis is aligned with $\vec{r}$ but
points in the opposite direction, toward the center of the earth.

The $x$ and $y$ axes lie in a plane perpendicular to $\vec{r}$ called the local
geocentric horizon plane. This plane is tangent to the spherical earth at the point
where $\vec{r}$ intersects the surface, point $S$ in
\cref{fig:local-geocentric-horizon}. The $x$ axis points toward the north pole, and
the $y$ axis points east, parallel to the equator.

\begin{figure}[htbp]
  \centering
  \resizebox{0.66\linewidth}{!}{\begin{tikzpicture}[scale=0.95,>=Latex,line cap=round,line join=round]
  \coordinate (O) at (0,0);
  \coordinate (S) at (2.55,2.05);
  \coordinate (V) at (2.93,2.36);
  % Earth section and equatorial plane.
  \draw[line width=0.9pt] (0,0) arc[start angle=180,end angle=45,radius=3.25];
  \draw[->,line width=0.9pt] (O) -- (4.95,0) node[right] {Equatorial plane};
  \draw[->,line width=0.9pt] (O) -- (0,4.20) node[above] {$\hat{z}_{\oplus}$};
  \node[align=center,left] at (-0.05,-0.05) {Earth\\center};
  % Position vector and vehicle.
  \draw[->,line width=1pt] (O) -- (V) node[midway,left] {$\vec r$};
  \fill (S) circle (1.5pt) node[below left] {$S$};
  \fill (V) circle (1.6pt);
  \node[right=6pt of V] {Vehicle};
  % Tangent horizon plane.
  \draw[line width=0.9pt] ($(V)+(-1.45,0.75)$) -- ($(V)+(1.45,-0.75)$);
  \node[align=left,right] at (4.15,1.55) {Local geocentric\\horizon plane};
  % Local axes: z inward, x north in tangent plane, y east out of section.
  \draw[->,line width=1pt] (V) -- ($(V)+(-0.90,0.47)$) node[above left] {$\hat{x}_{gc}$};
  \draw[->,line width=1pt] (V) -- ($(V)+(-0.70,-0.57)$) node[below left] {$\hat{z}_{gc}$};
  \draw[->,line width=1pt] (V) -- ($(V)+(0.15,0.88)$) node[above] {$\hat{y}_{gc}$};
\end{tikzpicture}
}
  \caption{Local geocentric horizon coordinates.}
  \label{fig:local-geocentric-horizon}
\end{figure}

The unit vectors, expressed in ECFC coordinates, are computed by
\texttt{geoc\_unit\_vectors}:
\begin{equation}
\begin{aligned}
  \hat{x}_{gc}
    &= -\sin\delta_{gc}
       (\cos\lambda\,\hat{x}_{\oplus}
        +\sin\lambda\,\hat{y}_{\oplus})
       +\cos\delta_{gc}\,\hat{z}_{\oplus},\\
  \hat{y}_{gc}
    &= -\sin\lambda\,\hat{x}_{\oplus}
       +\cos\lambda\,\hat{y}_{\oplus},\\
  \hat{z}_{gc}
    &= -\cos\delta_{gc}
       (\cos\lambda\,\hat{x}_{\oplus}
        +\sin\lambda\,\hat{y}_{\oplus})
       -\sin\delta_{gc}\,\hat{z}_{\oplus}.
\end{aligned}
\label{eq:local-geocentric-unit-vectors}
\end{equation}
Here $\delta_{gc}$ is geocentric latitude and $\lambda$ is longitude, defined in
\cref{sec:geocentric}.
